%!TeX root=../emmatop.tex
\chapter[Chapter \thechapter]{}
\lettrine[lines=4,lraise=0.3]{M}{r} Woodhouse was soon ready for his tea; and when he had drank his tea he was quite ready to go home; and it was as much as his three companions could do, to entertain away his notice of the lateness of the hour, before the other gentlemen appeared. Mr Weston was chatty and convivial, and no friend to early separations of any sort; but at last the drawing-room party did receive an augmentation. Mr Elton, in very good spirits, was one of the first to walk in. Mrs Weston and Emma were sitting together on a sofa. He joined them immediately, and, with scarcely an invitation, seated himself between them.

Emma, in good spirits too, from the amusement afforded her mind by the expectation of Mr Frank Churchill, was willing to forget his late improprieties, and be as well satisfied with him as before, and on his making Harriet his very first subject, was ready to listen with most friendly smiles.

He professed himself extremely anxious about her fair friend—her fair, lovely, amiable friend. »Did she know?—had she heard any thing about her, since their being at Randalls?—he felt much anxiety—he must confess that the nature of her complaint alarmed him considerably.« And in this style he talked on for some time very properly, not much attending to any answer, but altogether sufficiently awake to the terror of a bad sore throat; and Emma was quite in charity with him.

But at last there seemed a perverse turn; it seemed all at once as if he were more afraid of its being a bad sore throat on her account, than on Harriet's—more anxious that she should escape the infection, than that there should be no infection in the complaint. He began with great earnestness to entreat her to refrain from visiting the sick-chamber again, for the present—to entreat her to promise him not to venture into such hazard till he had seen Mr Perry and learnt his opinion; and though she tried to laugh it off and bring the subject back into its proper course, there was no putting an end to his extreme solicitude about her. She was vexed. It did appear—there was no concealing it—exactly like the pretence of being in love with her, instead of Harriet; an inconstancy, if real, the most contemptible and abominable! and she had difficulty in behaving with temper. He turned to Mrs Weston to implore her assistance, »Would not she give him her support?—would not she add her persuasions to his, to induce Miss Woodhouse not to go to Mrs Goddard's till it were certain that Miss Smith's disorder had no infection? He could not be satisfied without a promise—would not she give him her influence in procuring it?«

\begin{figure}[tbph]
\centering
\includegraphics[width=.9\linewidth]{15fair}
\caption{»Is this fair, Mrs Weston?«}
\end{figure}

»So scrupulous for others,« he continued, »and yet so careless for herself! She wanted me to nurse my cold by staying at home to-day, and yet will not promise to avoid the danger of catching an ulcerated sore throat herself. Is this fair, Mrs Weston?—Judge between us. Have not I some right to complain? I am sure of your kind support and aid.«

Emma saw Mrs Weston's surprize, and felt that it must be great, at an address which, in words and manner, was assuming to himself the right of first interest in her; and as for herself, she was too much provoked and offended to have the power of directly saying any thing to the purpose. She could only give him a look; but it was such a look as she thought must restore him to his senses, and then left the sofa, removing to a seat by her sister, and giving her all her attention.

She had not time to know how Mr Elton took the reproof, so rapidly did another subject succeed; for Mr John Knightley now came into the room from examining the weather, and opened on them all with the information of the ground being covered with snow, and of its still snowing fast, with a strong drifting wind; concluding with these words to Mr Woodhouse:

»This will prove a spirited beginning of your winter engagements, sir. Something new for your coachman and horses to be making their way through a storm of snow.«

Poor Mr Woodhouse was silent from consternation; but every body else had something to say; every body was either surprized or not surprized, and had some question to ask, or some comfort to offer. Mrs Weston and Emma tried earnestly to cheer him and turn his attention from his son-in-law, who was pursuing his triumph rather unfeelingly.

»I admired your resolution very much, sir,« said he, »in venturing out in such weather, for of course you saw there would be snow very soon. Every body must have seen the snow coming on. I admired your spirit; and I dare say we shall get home very well. Another hour or two's snow can hardly make the road impassable; and we are two carriages; if one is blown over in the bleak part of the common field there will be the other at hand. I dare say we shall be all safe at Hartfield before midnight.«

Mr Weston, with triumph of a different sort, was confessing that he had known it to be snowing some time, but had not said a word, lest it should make Mr Woodhouse uncomfortable, and be an excuse for his hurrying away. As to there being any quantity of snow fallen or likely to fall to impede their return, that was a mere joke; he was afraid they would find no difficulty. He wished the road might be impassable, that he might be able to keep them all at Randalls; and with the utmost good-will was sure that accommodation might be found for every body, calling on his wife to agree with him, that with a little contrivance, every body might be lodged, which she hardly knew how to do, from the consciousness of there being but two spare rooms in the house.

»What is to be done, my dear Emma?—what is to be done?« was Mr Woodhouse's first exclamation, and all that he could say for some time. To her he looked for comfort; and her assurances of safety, her representation of the excellence of the horses, and of James, and of their having so many friends about them, revived him a little.

His eldest daughter's alarm was equal to his own. The horror of being blocked up at Randalls, while her children were at Hartfield, was full in her imagination; and fancying the road to be now just passable for adventurous people, but in a state that admitted no delay, she was eager to have it settled, that her father and Emma should remain at Randalls, while she and her husband set forward instantly through all the possible accumulations of drifted snow that might impede them.

»You had better order the carriage directly, my love,« said she; »I dare say we shall be able to get along, if we set off directly; and if we do come to any thing very bad, I can get out and walk. I am not at all afraid. I should not mind walking half the way. I could change my shoes, you know, the moment I got home; and it is not the sort of thing that gives me cold.«

»Indeed!« replied he. »Then, my dear Isabella, it is the most extraordinary sort of thing in the world, for in general every thing does give you cold. Walk home!—you are prettily shod for walking home, I dare say. It will be bad enough for the horses.«

Isabella turned to Mrs Weston for her approbation of the plan. Mrs Weston could only approve. Isabella then went to Emma; but Emma could not so entirely give up the hope of their being all able to get away; and they were still discussing the point, when Mr Knightley, who had left the room immediately after his brother's first report of the snow, came back again, and told them that he had been out of doors to examine, and could answer for there not being the smallest difficulty in their getting home, whenever they liked it, either now or an hour hence. He had gone beyond the sweep—some way along the Highbury road—the snow was nowhere above half an inch deep—in many places hardly enough to whiten the ground; a very few flakes were falling at present, but the clouds were parting, and there was every appearance of its being soon over. He had seen the coachmen, and they both agreed with him in there being nothing to apprehend.

To Isabella, the relief of such tidings was very great, and they were scarcely less acceptable to Emma on her father's account, who was immediately set as much at ease on the subject as his nervous constitution allowed; but the alarm that had been raised could not be appeased so as to admit of any comfort for him while he continued at Randalls. He was satisfied of there being no present danger in returning home, but no assurances could convince him that it was safe to stay; and while the others were variously urging and recommending, Mr Knightley and Emma settled it in a few brief sentences: thus—

»Your father will not be easy; why do not you go?«

»I am ready, if the others are.«

»Shall I ring the bell?«

»Yes, do.«

And the bell was rung, and the carriages spoken for. A few minutes more, and Emma hoped to see one troublesome companion deposited in his own house, to get sober and cool, and the other recover his temper and happiness when this visit of hardship were over.

The carriage came: and Mr Woodhouse, always the first object on such occasions, was carefully attended to his own by Mr Knightley and Mr Weston; but not all that either could say could prevent some renewal of alarm at the sight of the snow which had actually fallen, and the discovery of a much darker night than he had been prepared for. »He was afraid they should have a very bad drive. He was afraid poor Isabella would not like it. And there would be poor Emma in the carriage behind. He did not know what they had best do. They must keep as much together as they could;« and James was talked to, and given a charge to go very slow and wait for the other carriage.

Isabella stept in after her father; John Knightley, forgetting that he did not belong to their party, stept in after his wife very naturally; so that Emma found, on being escorted and followed into the second carriage by Mr Elton, that the door was to be lawfully shut on them, and that they were to have a tête-à-tête drive. It would not have been the awkwardness of a moment, it would have been rather a pleasure, previous to the suspicions of this very day; she could have talked to him of Harriet, and the three-quarters of a mile would have seemed but one. But now, she would rather it had not happened. She believed he had been drinking too much of Mr Weston's good wine, and felt sure that he would want to be talking nonsense.

To restrain him as much as might be, by her own manners, she was immediately preparing to speak with exquisite calmness and gravity of the weather and the night; but scarcely had she begun, scarcely had they passed the sweep-gate and joined the other carriage, than she found her subject cut up—her hand seized—her attention demanded, and Mr Elton actually making violent love to her: availing himself of the precious opportunity, declaring sentiments which must be already well known, hoping—fearing—adoring—ready to die if she refused him; but flattering himself that his ardent attachment and unequalled love and unexampled passion could not fail of having some effect, and in short, very much resolved on being seriously accepted as soon as possible. It really was so. Without scruple—without apology—without much apparent diffidence, Mr Elton, the lover of Harriet, was professing himself her lover. She tried to stop him; but vainly; he would go on, and say it all. Angry as she was, the thought of the moment made her resolve to restrain herself when she did speak. She felt that half this folly must be drunkenness, and therefore could hope that it might belong only to the passing hour. Accordingly, with a mixture of the serious and the playful, which she hoped would best suit his half and half state, she replied,

»I am very much astonished, Mr Elton. This to me! you forget yourself—you take me for my friend—any message to Miss Smith I shall be happy to deliver; but no more of this to me, if you please.«

»Miss Smith!—message to Miss Smith!—What could she possibly mean!«—And he repeated her words with such assurance of accent, such boastful pretence of amazement, that she could not help replying with quickness,

»Mr Elton, this is the most extraordinary conduct! and I can account for it only in one way; you are not yourself, or you could not speak either to me, or of Harriet, in such a manner. Command yourself enough to say no more, and I will endeavour to forget it.«

But Mr Elton had only drunk wine enough to elevate his spirits, not at all to confuse his intellects. He perfectly knew his own meaning; and having warmly protested against her suspicion as most injurious, and slightly touched upon his respect for Miss Smith as her friend,—but acknowledging his wonder that Miss Smith should be mentioned at all,—he resumed the subject of his own passion, and was very urgent for a favourable answer.

As she thought less of his inebriety, she thought more of his inconstancy and presumption; and with fewer struggles for politeness, replied,

»It is impossible for me to doubt any longer. You have made yourself too clear. Mr Elton, my astonishment is much beyond any thing I can express. After such behaviour, as I have witnessed during the last month, to Miss Smith—such attentions as I have been in the daily habit of observing—to be addressing me in this manner—this is an unsteadiness of character, indeed, which I had not supposed possible! Believe me, sir, I am far, very far, from gratified in being the object of such professions.«

»Good Heaven!« cried Mr Elton, »what can be the meaning of this?—Miss Smith!—I never thought of Miss Smith in the whole course of my existence—never paid her any attentions, but as your friend: never cared whether she were dead or alive, but as your friend. If she has fancied otherwise, her own wishes have misled her, and I am very sorry—extremely sorry—But, Miss Smith, indeed!—Oh! Miss Woodhouse! who can think of Miss Smith, when Miss Woodhouse is near! No, upon my honour, there is no unsteadiness of character. I have thought only of you. I protest against having paid the smallest attention to any one else. Every thing that I have said or done, for many weeks past, has been with the sole view of marking my adoration of yourself. You cannot really, seriously, doubt it. No!—(in an accent meant to be insinuating)—I am sure you have seen and understood me.«

It would be impossible to say what Emma felt, on hearing this—which of all her unpleasant sensations was uppermost. She was too completely overpowered to be immediately able to reply: and two moments of silence being ample encouragement for Mr Elton's sanguine state of mind, he tried to take her hand again, as he joyously exclaimed—

»Charming Miss Woodhouse! allow me to interpret this interesting silence. It confesses that you have long understood me.«

»No, sir,« cried Emma, »it confesses no such thing. So far from having long understood you, I have been in a most complete error with respect to your views, till this moment. As to myself, I am very sorry that you should have been giving way to any feelings—Nothing could be farther from my wishes—your attachment to my friend Harriet—your pursuit of her, (pursuit, it appeared,) gave me great pleasure, and I have been very earnestly wishing you success: but had I supposed that she were not your attraction to Hartfield, I should certainly have thought you judged ill in making your visits so frequent. Am I to believe that you have never sought to recommend yourself particularly to Miss Smith?—that you have never thought seriously of her?«

»Never, madam,« cried he, affronted in his turn: »never, I assure you. I think seriously of Miss Smith!—Miss Smith is a very good sort of girl; and I should be happy to see her respectably settled. I wish her extremely well: and, no doubt, there are men who might not object to—Every body has their level: but as for myself, I am not, I think, quite so much at a loss. I need not so totally despair of an equal alliance, as to be addressing myself to Miss Smith!—No, madam, my visits to Hartfield have been for yourself only; and the encouragement I received\longdash«

»Encouragement!—I give you encouragement!—Sir, you have been entirely mistaken in supposing it. I have seen you only as the admirer of my friend. In no other light could you have been more to me than a common acquaintance. I am exceedingly sorry: but it is well that the mistake ends where it does. Had the same behaviour continued, Miss Smith might have been led into a misconception of your views; not being aware, probably, any more than myself, of the very great inequality which you are so sensible of. But, as it is, the disappointment is single, and, I trust, will not be lasting. I have no thoughts of matrimony at present.«

He was too angry to say another word; her manner too decided to invite supplication; and in this state of swelling resentment, and mutually deep mortification, they had to continue together a few minutes longer, for the fears of Mr Woodhouse had confined them to a foot-pace. If there had not been so much anger, there would have been desperate awkwardness; but their straightforward emotions left no room for the little zigzags of embarrassment. Without knowing when the carriage turned into Vicarage Lane, or when it stopped, they found themselves, all at once, at the door of his house; and he was out before another syllable passed.—Emma then felt it indispensable to wish him a good night. The compliment was just returned, coldly and proudly; and, under indescribable irritation of spirits, she was then conveyed to Hartfield.

There she was welcomed, with the utmost delight, by her father, who had been trembling for the dangers of a solitary drive from Vicarage Lane—turning a corner which he could never bear to think of—and in strange hands—a mere common coachman—no James; and there it seemed as if her return only were wanted to make every thing go well: for Mr John Knightley, ashamed of his ill-humour, was now all kindness and attention; and so particularly solicitous for the comfort of her father, as to seem—if not quite ready to join him in a basin of gruel—perfectly sensible of its being exceedingly wholesome; and the day was concluding in peace and comfort to all their little party, except herself.—But her mind had never been in such perturbation; and it needed a very strong effort to appear attentive and cheerful till the usual hour of separating allowed her the relief of quiet reflection.